\chapter{Complexity — Logspace Tableau}
%%% generated by the blueprint pipeline from TCSlib.Complexity.CircuitComplexity.LogspaceTableau
\section{Overview}

Declarations of the Lean module \texttt{TCSlib.Complexity.CircuitComplexity.LogspaceTableau}.

\section{Declarations}

\begin{theorem}[Front words have logarithmic-space query languages]
\lean{Complexity.unaryLogspace_frontWord}
\label{Complexity.unaryLogspace_frontWord}
\leanok
\uses{Complexity.UTab.Front.bnd, Complexity.UTab.Front.prog, Complexity.UTab.Front.run_word, Complexity.UTab.Front.word, Complexity.UnaryLogspace, Complexity.UnaryLogspace.counterProg, Complexity.unaryLogspace_replicate}
\difficulty{1}
For all $C,d\in\mathbb N$, define $w(n)$ to be $n$ ones, followed by $0$, then
$(C+1)(n+1)^d$ ones, followed by $0$. Both languages of encoded pairs consisting of
$1^n$ and the binary representation of an index—one accepting when the indexed bit of
$w(n)$ is $1$, the other when the index is less than $|w(n)|$—are decidable in
logarithmic space.
\end{theorem}

\begin{theorem}[Bounded runtime of the circuit encoder]
\lean{Complexity.emit_run}
\label{Complexity.emit_run}
\leanok
\uses{BoolCircuit.DAGCircuit.encode, Complexity.CounterProg.init, Complexity.CounterProg.run, Complexity.UTab.Front.tb, Complexity.UTab.Front.word, Complexity.UTab.emitC, Complexity.UTab.goes_any, Complexity.UTab.prog, Complexity.UTab.run_budget, Complexity.UTab.tabEmit, Complexity.UTab.tabEmit_eq, Complexity.tabFamily}
\difficulty{5}
Let $M$ be a finite Turing machine with Boolean input and output symbols, and let
$C,d,n$ be natural numbers. There is a time $t \le 60(k+1)(L+30)(C+4)^3(n+1)^{3(d+1)}$,
where $k$ is the number of work tapes of $M$ and $L$ is the maximum length of the
instruction templates in the counter program constructed from $M$, such that this
program, started with all eleven registers equal to zero, at its initial label, and with
empty output, halts after $t$ steps on the word $1^n0\,1^{(C+1)(n+1)^d}0$ and has output
equal to the encoding of the circuit at input length $n$ in the family constructed from
$M,C,d$.
\end{theorem}

\begin{theorem}[Logspace uniformity of the tableau family]
\lean{Complexity.tabFamily_isLogspaceUniform}
\label{Complexity.tabFamily_isLogspaceUniform}
\leanok
\uses{BoolCircuit.DAGCircuit.encode, BoolCircuit.DAGCircuitFamily.IsLogspaceUniform, Complexity.CounterProg.length_out_le, Complexity.UTab.emitC, Complexity.UTab.prog, Complexity.UnaryLogspace, Complexity.UnaryLogspace.counterProg, Complexity.UnaryLogspace.implicitlyLogspaceComputable, Complexity.emit_run, Complexity.tabFamily, Complexity.unaryExt, Complexity.unaryExt_replicate, Complexity.unaryLogspace_frontWord}
\difficulty{4}
For every finite-state Turing machine $M$ with Boolean input and output alphabets and
all natural numbers $C,d$, the circuit family whose $n$th circuit is the
configuration-tableau circuit for $M$ with time bound $(C+1)(n+1)^d$ is
logspace-uniform.
\end{theorem}

\begin{theorem}[Logspace-uniform circuits for polynomial-time languages]
\lean{Language.exists_isLogspaceUniform_of_mem_P}
\label{Language.exists_isLogspaceUniform_of_mem_P}
\leanok
\uses{BoolCircuit.DAGCircuitFamily, BoolCircuit.DAGCircuitFamily.HasFaninTwo, BoolCircuit.DAGCircuitFamily.IsLogspaceUniform, BoolCircuit.DAGCircuitFamily.IsLogspaceUniform.isPolySize, BoolCircuit.DAGCircuitFamily.IsPolySize, BoolCircuit.DAGCircuitFamily.language, Complexity.P, Complexity.mem_P_iff, Complexity.tabFamily, Complexity.tabFamily_hasFaninTwo, Complexity.tabFamily_isLogspaceUniform, Complexity.tabFamily_language}
\difficulty{1}
Every language of Boolean strings decidable in polynomial time has a logspace-uniform
family of Boolean circuits with at most two inputs per gate, of polynomial size, whose
accepted language is exactly the given language.
\end{theorem}

\begin{theorem}[Logspace-uniform circuit characterization of P]
\lean{Language.mem_P_iff_exists_isLogspaceUniform}
\label{Language.mem_P_iff_exists_isLogspaceUniform}
\leanok
\uses{BoolCircuit.DAGCircuitFamily, BoolCircuit.DAGCircuitFamily.HasFaninTwo, BoolCircuit.DAGCircuitFamily.IsLogspaceUniform, BoolCircuit.DAGCircuitFamily.IsLogspaceUniform.isPUniform, BoolCircuit.DAGCircuitFamily.IsPolySize, BoolCircuit.DAGCircuitFamily.language, Complexity.P, Language.exists_isLogspaceUniform_of_mem_P, Language.mem_P_of_isPUniform}
\difficulty{1}
A language $L$ over the Boolean alphabet belongs to polynomial time if and only if there
is a family of Boolean circuits that recognizes exactly $L$, has fan-in at most two, has
size bounded by a polynomial in the input length, and is logspace-uniform.
\end{theorem}
